\documentclass[a4paper, serif, 10pt]{moderncv}

%% moderncv
% style and color
\moderncvstyle{banking}
\moderncvcolor{black}

%% page layout/margins
\usepackage[top=2.5cm, bottom=2.5cm, left=1.5cm, right=1.5cm]{geometry}

\nopagenumbers{}

%% Babel config for Dutch language
% ref:
% - https://latex3.github.io/babel/guides/locale-dutch.html
\usepackage[dutch]{babel}

%% fontspec
\usepackage{fontspec}

\IfFontExistsTF{Linux Libertine}{
  \setmainfont[Ligatures={TeX, Common}, Numbers=OldStyle]{Linux Libertine}
}{}
\IfFontExistsTF{Linux Libertine O}{
  \setmainfont[Ligatures={TeX, Common}, Numbers=OldStyle]{Linux Libertine O}
}{}

%% CTeX
% CJK support, used to show CJK characters in the resume
%
% - fontset=none: disable builtin fontset but instead set the CJK font manually
% - heading=false: disable ctex heading
% - punct=kaiming: use kaiming punctuations styles for CJK
% - scheme=plain: use plain scheme, do not override \normalsize font size
% - space=auto: space settings for CJK characters
%
% ref:
% - http://ctan.mirrorcatalogs.com/language/chinese/ctex/ctex.pdf
\usepackage[UTF8, heading=false, punct=kaiming, scheme=plain, space=auto]{ctex}

\IfFontExistsTF{Noto Serif CJK SC}{
  \setCJKmainfont{Noto Serif CJK SC}
}{}
\IfFontExistsTF{Noto Sans CJK SC}{
  \setCJKsansfont{Noto Sans CJK SC}
}{}

%% line spacing
\usepackage{setspace}
\setstretch{1.125}

%% URL styling - use normal text instead of monospace
\urlstyle{same}

% Auto-underline all links
% Defer until \begin{document} so hyperref is already loaded
\AtBeginDocument{%
  \let\oldhref\href
  \renewcommand{\href}[2]{\oldhref{#1}{\underline{#2}}}%
}

%% Patch \httplink and \httpslink to handle full URLs
% The original moderncv macros always prepend http:// or https:// to URLs.
% This causes issues when the URL already has a protocol (e.g., https://example.com)
% resulting in malformed URLs like https://https://example.com.
% This patch checks if the URL already contains "://" and uses it directly if so.
% Note: We use \str_set:Nx (x-expansion) to expand macros like \@homepage first.
\ExplSyntaxOn
\str_new:N \g_yamlresume_url_str

\RenewDocumentCommand{\httpslink}{O{}m}{
  \str_set:Nx \g_yamlresume_url_str {#2}
  \str_if_in:NnTF \g_yamlresume_url_str {://}
    { \url{#2} }
    { \url{https://#2} }
}

\RenewDocumentCommand{\httplink}{O{}m}{
  \str_set:Nx \g_yamlresume_url_str {#2}
  \str_if_in:NnTF \g_yamlresume_url_str {://}
    { \url{#2} }
    { \url{http://#2} }
}
\ExplSyntaxOff

\name{Sanne van den Berg}{}
\title{Senior UX-designer met focus op toegankelijk en inclusief ontwerp}
\phone[mobile]{+31 6 12345678}
\email{sanne@vandenberg.design}
\homepage{https://sannevandenberg.design}

\address{Keizersgracht 123, Amsterdam, Noord-Holland, Nederland, 1015 CJ}{}{}

\extrainfo{{\small \faLinkedin}\ \href{https://www.linkedin.com/in/sannevandenberg}{@sannevandenberg} {} {} {} • {} {} {}
{\small \faBehance}\ \href{https://www.behance.net/sannevandenberg}{@sannevandenberg} {} {} {} • {} {} {}
{\small \faDribbble}\ \href{https://dribbble.com/sannevandenberg}{@sannevandenberg}}

\begin{document}

\maketitle

\section{Basis Informatie}

\cvline{}{Ontwerper met ruim 8 jaar ervaring in het creëren van gebruiksvriendelijke
digitale producten voor webshops, SaaS-toepassingen en publieke diensten.
Ik vertaal gebruikersbehoeften naar heldere interacties en inclusieve
ontwerpen, met oog voor zowel businessdoelen als menselijke waarden.
%
\begin{itemize}
\item Methodisch en onderzoekend: van contextanalyse tot usabilitytest
\item Ervaren met design systems, rapid prototyping en service design
\item Gedreven om digitale toegankelijkheid (WCAG 2.2) standaard te maken
\end{itemize}}
\section{Opleidingen}

\cventry{sep 2010–jul 2014}
        {Bachelor, Communicatie en Multimedia Design, Score: 8.2}
        {Hogeschool van Amsterdam}
        {\url{https://www.hva.nl}}
        {}
        {\begin{itemize}
\item Afstudeerrichting Interaction Design, afgestudeerd met een 8,2
\item Afstudeerproject: toegankelijke tabletapp voor senioren
\item Minor: Design for Social Change
\end{itemize}
\textbf{Cursussen}: Onderzoeksmethoden voor Design, Interactieontwerp, Visual Interface Design, Mens-computerinteractie, Digitale Productie, Concepting en Creativiteit, Bedrijfskunde voor Creatieve Industrie, Afstudeerproject Apps voor Ouderen}
\section{Werkervaring}

\cventry{mrt 2021–Heden}
        {Senior UX Designer}
        {Bluebird Digital}
        {\url{https://www.bluebird-digital.nl}}
        {}
        {\begin{itemize}
\item Verantwoordelijk voor de UX-strategie van een B2B SaaS-platform met ruim 40.000 gebruikers.
\item Opgezet en doorontwikkeld design system met toegankelijke, herbruikbare componenten in Figma.
\item Begeleidde junior designers en gaf UX-workshops aan developers en productmanagers.
\end{itemize}
\textbf{Trefwoorden}: Design system, SaaS, WCAG 2.2, Agile/Scrum, Mentorschap}

\cventry{mei 2017–feb 2021}
        {Product Designer}
        {Startup Studio Amsterdam}
        {\url{https://www.startupstudio.amsterdam}}
        {}
        {\begin{itemize}
\item Ontwikkelde UX-ontwerpen voor 12 startups in fintech, health en mobility.
\item Leidde wekelijks usabilitytesten en vertaalde inzichten naar prototypes.
\item Creëerde customer journeys en service blueprints voor complexe processen.
\end{itemize}
\textbf{Trefwoorden}: Fintech, Mobility, Usabilitytesten, Customer journey, Prototyping}

\cventry{aug 2014–apr 2017}
        {Junior UX Designer}
        {Glow Online}
        {\url{https://www.glowonline.nl}}
        {}
        {\begin{itemize}
\item Ontwierp wireframes en interactieve prototypes voor e-commerce websites.
\item Voerde A/B-tests uit om conversie en gebruiksvriendelijkheid te verbeteren.
\item Werkte nauw samen met developers bij de implementatie van responsive designs.
\end{itemize}
\textbf{Trefwoorden}: E-commerce, Wireframes, A/B-testen, Responsive design, Conversie}
\section{Talen}

\cvline{Nederlands}{Moedertaal of tweetalige vaardigheid \hfill \textbf{Trefwoorden}: Moedertaal}
\cvline{Engels}{Volledige professionele vaardigheid \hfill \textbf{Trefwoorden}: Professioneel C1}
\cvline{Duits}{Beperkte werkvaardigheid \hfill \textbf{Trefwoorden}: Basisvaardigheden}
\section{Vaardigheden}

\cvline{User Research}{Expert \hfill \textbf{Trefwoorden}: Diepte-interviews, Usabilitytesten, A/B-testen, Contextanalyse, Persona's}
\cvline{Prototyping}{Expert \hfill \textbf{Trefwoorden}: Figma, Interactieve prototypes, Wireframes, Design tokens, Clickdummy's}
\cvline{Interaction Design}{Geavanceerd \hfill \textbf{Trefwoorden}: Customer journeys, Service blueprints, Micro-interacties, Toegankelijk ontwerp}
\cvline{Visual Design}{Geavanceerd \hfill \textbf{Trefwoorden}: Design systems, Typografie, Kleurtoegankelijkheid, WCAG 2.2}
\cvline{Front-end Basics}{Intermediair \hfill \textbf{Trefwoorden}: HTML, CSS, JavaScript, Storybook}
\section{Onderscheidingen}

\cventry{nov 2022}
        {Dutch UX Association}
        {Publieksprijs Nederlandse UX Awards}
        {}
        {}
        {Bekroond voor het toegankelijk herontwerpen van een overheidsportaal voor
mensen met een visuele beperking.}
\section{Certificaten}

\cventry{jun 2023}
        {IAAP}
        {Certified Professional in Accessibility Core Competencies (CPACC)}
        {\url{https://www.accessibilityassociation.org/}}
        {}
        {}

\cventry{apr 2021}
        {Coursera}
        {Google UX Design Professional Certificate}
        {\url{https://www.coursera.org/professional-certificates/google-ux-design}}
        {}
        {}
\section{Publicaties}

\cventry{mei 2023}
        {Inclusief ontwerpen is geen trend, het is een vak}
        {UX Magazine Nederland}
        {\url{https://www.uxmagazine.nl/artikel/inclusief-ontwerpen}}
        {}
        {Praktijkartikel over hoe je toegankelijkheid vanaf de start van het
ontwerpproces meeneemt, met concrete voorbeelden uit de financiële sector.}
\section{Referenties}

\cventry{\emaillink[liesbeth.jansen@bluebird-digital.nl]{liesbeth.jansen@bluebird-digital.nl}}
        {Design lead bij Bluebird Digital}
        {Dr. Liesbeth Jansen}
        {+31 20 555 0199}
        {}
        {Sanne is een van de meest zorgvuldige en empathische designers met wie ik
heb gewerkt. Haar aanpak combineert onderzoek, creativiteit en oog voor
detail met een sterke drive om digitale producten voor iedereen bruikbaar te maken.}
\section{Projecten}

\cventry{jan 2023–jun 2023}
        {Volledig herontwerp van de mobiele bankapp voor een grote Nederlandse
bank, gericht op begrijpelijke taal, donkere modus en volledige
ondersteuning van schermlezers.
}
        {OmniBank App Herontwerp}
        {\url{https://sannevandenberg.design/projecten/omnibank}}
        {}
        {\begin{itemize}
\item Uitgebreid gebruikersonderzoek met 18 deelnemers, inclusief mensen met een visuele beperking.
\item Iteratieve prototypes getest in vijf rondes op toegankelijkheid.
\item Opgeleverd design package met componenten, tokens en richtlijnen voor developers.
\end{itemize}
\textbf{Trefwoorden}: Fintech, Toegankelijkheid, Figma, Co-creatie, Design system}
\section{Interesses}

\cvline{Duurzaamheid}{Circulair design, Tweedehands meubels, Stadsmoestuin}
\cvline{Fotografie}{Straatfotografie, Analoog, Zwart-wit}
\cvline{Muziek}{Gitaar, Indie, Vinyl}
\section{Vrijwilligerswerk}

\cventry{jan 2019–Heden}
        {Vrijwillig UX-adviseur}
        {Inclusive Design Collectief}
        {\url{https://www.inclusiefdesigncollectief.nl}}
        {}
        {\begin{itemize}
\item Help non-profitorganisaties met het toegankelijk maken van hun digitale kanalen.
\item Organiseer jaarlijkse toegankelijkheidssessies tijdens de Dutch Design Week.
\end{itemize}}

\end{document}